\subsubsection*{命题 4.3.}
\underline{设 $x_1$ 和 $x_2$ 为 $\mathbb P^r$ 的两个不同闭点, $L_1$ 和 $L_2$ 为 $\mathbb P^r$ 的两个线性子空间, 且 $x_1\in L_1$, $x_2\in L_2$. 在 $d$ 次超曲面的空间}
\[
\check{\mathbb P}^{N}
=\check{\mathbb P}^{\binom{r+d}{d}-1}
\]
\underline{中, 由在 $x_1$ 处与 $L_1$ 相切且在 $x_2$ 处与 $L_2$ 相切的超曲面所成的线性子空间的余维为}
\[
\left\{
\begin{aligned}
\dim L_1+\dim L_2+2,
 &\quad\text{若 $d\geq3$, 或 $d=2$ 且 $L_1\cap L_2$ 不}\\[-0.2ex]
 &\quad\text{包含经过 $x_1,x_2$ 的直线 $(x_1,x_2)$},\\[1ex]
\dim L_1+\dim L_2+1,
 &\quad\text{若 $d=2$ 且 $L_1\cap L_2$ 包含直线 $(x_1,x_2)$.}
\end{aligned}
\right.
\]

\noindent\underline{证明.}
对 $d\geq3$, 须证明 $\dim L_1+\dim L_2+2$ 个线性条件彼此独立. 因而只需考虑 $L_1=L_2=\mathbb P^r$ 的情形, 即确定使一个超曲面在 $x_1$ 和 $x_2$ 处都有奇点需要多少个独立条件. 选取齐次坐标, 使
\[
x_1=(1,0,\ldots,0),\qquad x_2=(0,1,0,\ldots,0).
\]
由方程 $\sum_W A_WX^W=0$ 定义的超曲面在 $(1,0,\ldots,0)$ 处有奇点, 当且仅当所有单项式 $X_0^{d-1}X_i$, $i=0,\ldots,r$, 的系数均为零. 它在 $(0,1,0,\ldots,0)$ 处有奇点, 当且仅当所有单项式 $X_1^{d-1}X_i$, $i=0,\ldots,r$, 的系数均为零. 当 $d\geq3$ 时这些单项式互不相同, 这就证明了此情形的断言.

对 $d=2$, 同样的计算给出
\[
\operatorname{codim}\geq\dim L_1+\dim L_2+1,
\]
因为单项式 $X_0X_1$ 被计算了两次.

先假设 $L_1\cap L_2$ 不包含直线 $(x_1,x_2)$. 归结到 $L_1=\mathbb P^r$ 且 $L_2$ 为 $\mathbb P^r$ 中不包含 $(x_1,x_2)$ 的超平面的情形. 选齐次坐标 $X_0,\ldots,X_r$, 使
\[
x_1=(1,0,\ldots,0),\qquad x_2=(0,1,0,\ldots,0),
\qquad L_2=(X_0=0).
\]
二次超曲面 $\sum_{i\leq j}A_{ij}X_iX_j$ 在 $(1,0,\ldots,0)$ 处有奇点, 当且仅当对 $j=0,\ldots,r$ 有 $A_{0j}=0$; 它在 $(0,1,\ldots,0)$ 处与 $X_0=0$ 相切, 当且仅当对 $j=1,\ldots,r$ 有 $A_{1j}=0$. 合在一起共有
\[
2r+1=\dim L_1+\dim L_2+2
\]
个独立条件.

最后, 假设 $L_1\cap L_2$ 包含直线 $(x_1,x_2)$. 归结到 $L_1=L_2=(x_1,x_2)$ 的情形. 在满足
\[
x_1=(1,0,\ldots,0),\qquad x_2=(0,1,0,\ldots,0),
\]
的齐次坐标中, 二次超曲面 $\sum_{i\leq j}A_{ij}X_iX_j$ 经过 $(1,0,\ldots,0)$ 且在该处与 $L_1$ 相切, 当且仅当 $A_{00}=A_{01}=0$. 它经过 $(0,1,0,\ldots,0)$ 且在该处与 $L_2$ 相切, 当且仅当 $A_{11}=A_{01}=0$. 因此总共只有
\[
3=\dim L_1+\dim L_2+1
\]
个条件.